% Generated draft: requires mathematical and source review.
\documentclass[11pt]{article}
\usepackage[margin=25mm]{geometry}
\usepackage{fontspec}
\setmainfont{Latin Modern Roman}
\usepackage{amsmath,amssymb,mathtools,unicode-math}
\setmathfont{Latin Modern Math}
\usepackage{enumitem,needspace,xurl,hyperref}
\setlength{\parindent}{0pt}
\setlength{\parskip}{5pt}
\setlength{\emergencystretch}{3em}
\newcommand{\N}{\mathbb{N}}
\newcommand{\Z}{\mathbb{Z}}
\newcommand{\Q}{\mathbb{Q}}
\newcommand{\R}{\mathbb{R}}
\newcommand{\C}{\mathbb{C}}
\newcommand{\F}{\mathbb{F}}
\newcommand{\A}{\mathbb{A}}
\newcommand{\PP}{\mathbb{P}}
\newcommand{\E}{\mathbb{E}}
\newcommand{\eps}{\varepsilon}
\newcommand{\dd}{\,\mathrm{d}}
\begin{document}
\section*{UnsolvedMath ID: 3100072}
\textbf{Source problem code:} catalogue entry
\subsection{Definitions and mathematical statement}
Let $\mathbb F_2=\mathbb Z/2\mathbb Z$, and let
\[
f(X)=\sum_{n\geq0}a_nX^n\in\mathbb F_2[[X]],\qquad a_0=1.
\]
For $n\geq1$, suppose the coefficients $a_n$ are independent Bernoulli random variables with $\Pr(a_n=1)=p_n$ and $\Pr(a_n=0)=1-p_n$. The source describes $p_n$ using the expression ``\textnormal{Ü}$(n^{-2})$'', whose leading symbol is not a defined probability or asymptotic operator in the supplied record. Thus the sequence $(p_n)$, including its admissible quantifiers, remains unspecified.

Since $a_0=1$, there is a unique formal reciprocal
\[
\frac1{f(X)}=\sum_{n\geq0}b_nX^n\in\mathbb F_2[[X]],
\qquad b_0=1,\qquad
b_n=\sum_{k=1}^n a_k b_{n-k}\quad(n\geq1),
\]
where the last sum is in $\mathbb F_2$. Regard $b_n$ also as the integer $0$ or $1$. The source asks whether the density of ones in this reciprocal is $1/2$, namely whether
\[
\lim_{N\to\infty}\frac1N\sum_{n=0}^{N-1}b_n=\frac12.
\]
If an almost-sure assertion is intended, the precise probabilistic target is that this limit exists and equals $1/2$ with probability one. The source does not explicitly specify almost-sure convergence, convergence in probability, or an assertion about expected densities, and the corrupted probability expression prevents fixing a complete probability law.
\subsection{Short English statement}
Choose independent binary coefficients with constant term one and study the formal reciprocal over $\mathbb F_2$. Must half its coefficients be ones in the limiting natural-density sense? The probability formula is corrupted and the mode of probabilistic assertion is unspecified.
\subsection{Sources}
\begin{itemize}
\item Cooper, \emph{Combinatorial Problems I Like}, October 2020 snapshot, ``Classical Number Theory,'' source-order bullet 72, \url{https://people.math.sc.edu/cooper/combprob.html}; accessed 2026-07-29, HTML extraction.
\item The supplied research summary calls this a companion to source item 0071 associated with Riasanovsky; it supplies no further bibliographic citation or verified resolution.
\item Local registry \texttt{data/unsolvedmath/problems.json}, permanent UnsolvedMath ID $3100072$, source code catalogue entry. Dataset attribution: Ulam AI and the UnsolvedMath Contributors.
\end{itemize}
\end{document}
